\section{Introduction}
Reinforcement learning systems are often trained against a proxy reward while
being evaluated, implicitly or explicitly, against a broader intended
objective. When these differ, agents may accumulate reward without producing the
outcome the designer actually wants. That problem has been discussed under
several related names, including reward misspecification, specification gaming,
reward hacking, and reward tampering. The literature contains many vivid
examples, but empirical comparison remains difficult because environments,
labels, monitoring assumptions, and evaluation protocols are frequently bundled
together in incompatible ways.

This paper takes a framework-first position. Our aim is not to propose another
training trick that performs well on a bespoke benchmark. Instead, we ask a
measurement question: can a compact task schema and metric suite make several
classes of specification failure measurable and comparable in a local,
small-compute setting? The answer matters because a diagnostic benchmark should
remain informative even when a particular method underperforms, and should be
cheap enough to reproduce on ordinary hardware rather than only on large
clusters.

We formalize a specification diagnostic task as
\[
T = (M, R_p, J_t, C, M_{\mathrm{mon}}, \Delta, \Pi),
\]
where $M$ is the underlying MDP or POMDP, $R_p$ is the proxy reward used for
learning, $J_t$ is the hidden true objective used for evaluation, $C$ is the
corruption or tampering channel family, $M_{\mathrm{mon}}$ is the monitor
observation channel, $\Delta$ is the distribution-shift family, and $\Pi$ is
the policy or training family under consideration. This decomposition is meant
to make explicit what is often implicit: the environment dynamics, training
signal, evaluation target, monitoring assumptions, and robustness conditions are
distinct design choices and should be measured as such.

\framework packages that decomposition into code and paper artifacts. The
repository provides three symbolic environments, a schema export for each task,
a metric suite for quantifying specification gap and monitorability, leakage
diagnostics for audit channels, and a staged execution pipeline that respects
laptop-scale compute and disk budgets. The design goal is to support empirical
clarity rather than breadth. Symbolic environments are not intended as realistic
simulators of deployment settings; they are compact model organisms for studying
failure modes under controlled corruption and shift.

The empirical sections of this paper are intentionally constrained by a strict
artifact policy. Tables, plots, and result sentences are only populated from
generated files saved by the repository. If a run fails, a metric degenerates,
or leakage forces a feature redesign, the paper records that outcome rather than
smoothing it away. This makes the benchmark less polished than a hand-curated
final paper at early stages, but more honest as a research artifact.

Our main contribution is therefore a reproducible evaluation object rather than
an algorithmic claim: a task schema, metric protocol, benchmark instantiation,
and monitor evaluation procedure for studying reward misspecification and
tampering in reinforcement learning under modest hardware constraints.

Empirically, the current benchmark supports a conservative but useful claim:
specification gaps are easy to measure across all three environments, tampering
risk is directly measurable in SensorTamper-v0, and leakage-controlled monitor
evaluation is feasible but currently narrow, with accepted monitor evidence
concentrated in SensorTamper-v0.
